
\chapter{2016年北京卷（文）}

\section{单选题}

\begin{problem}
已知集合 \(A=\{x\mid 2<x<4\}\)，\(B=\{x\mid x<3 \text{或} x>5\}\)，则 \(A\cap B=\)（\quad）

\choices
    {\(\{x\mid 2<x<5\}\)}
    {\(\{x\mid x<4 \text{或} x>5\}\)}
    {\(\{x\mid 2<x<3\}\)}
    {\(\{x\mid x<2 \text{或} x>5\}\)}
\end{problem}

\begin{problem}
复数 \(\dfrac{1+2\i}{2-\i}=\)（\quad）

\choices
    {\(\i\)}
    {\(1+\i\)}
    {\(-\i\)}
    {\(1-\i\)}
\end{problem}

\begin{problem}
执行如图所示的程序框图，输出 \(s\) 的值为（\quad）

\bitmapfigure[max width=6.0cm]{img/2016/beijing_liberal/q3_fig1.png}

\choices
    {\(8\)}
    {\(9\)}
    {\(27\)}
    {\(36\)}
\end{problem}

\begin{problem}
下列函数中，在区间 \((-1,1)\) 上为减函数的是（\quad）

\choices
    {\(y=\dfrac{1}{1-x}\)}
    {\(y=\cos x\)}
    {\(y=\ln\left( x+1 \right)\)}
    {\(y=2^{-x}\)}
\end{problem}

\begin{problem}
圆 \(\left( x+1 \right)^2+y^2=2\) 的圆心到直线 \(y=x+3\) 的距离为（\quad）

\choices
    {\(1\)}
    {\(2\)}
    {\(\sqrt2\)}
    {\(2\sqrt2\)}
\end{problem}

\begin{problem}
从甲，乙等 \(5\) 名学生中随机选出 \(2\) 人，则甲被选中的概率为（\quad）

\choices
    {\(\dfrac15\)}
    {\(\dfrac25\)}
    {\(\dfrac{8}{25}\)}
    {\(\dfrac{9}{25}\)}
\end{problem}

\begin{problem}
已知 \(A(2,5)\)，\(B(4,1)\). 若点 \(P(x,y)\) 在线段 \(AB\) 上，则 \(2x-y\) 的最大值为（\quad）

\choices
    {\(-1\)}
    {\(3\)}
    {\(7\)}
    {\(8\)}
\end{problem}

\begin{problem}
某学校运动会的立定跳远和 \(30\) 秒跳绳两个单项比赛分成预赛和决赛两个阶段，表中为 \(10\) 名学生的预赛成绩，其中有三个数据模糊.

\begin{center}
\begin{tabular}{|c|c|c|c|c|c|c|c|c|c|c|}
\hline
学生序号 & 1 & 2 & 3 & 4 & 5 & 6 & 7 & 8 & 9 & 10 \\
\hline
立定跳远（单位：米） & 1.96 & 1.92 & 1.82 & 1.80 & 1.78 & 1.76 & 1.74 & 1.72 & 1.68 & 1.60 \\
\hline
\(30\) 秒跳绳（单位：次） & 63 & \(a\) & 75 & 60 & 63 & 72 & 70 & \(a-1\) & \(b\) & 65 \\
\hline
\end{tabular}
\end{center}

在这 \(10\) 名学生中，进入立定跳远决赛的有 \(8\) 人，同时进入立定跳远决赛和 \(30\) 秒跳绳决赛的有 \(6\) 人，则（\quad）

\choices
    {2号学生进入 \(30\) 秒跳绳决赛}
    {5号学生进入 \(30\) 秒跳绳决赛}
    {8号学生进入 \(30\) 秒跳绳决赛}
    {9号学生进入 \(30\) 秒跳绳决赛}
\end{problem}

\section{填空题}

\begin{problem}
已知向量 \(\bs a=(1,\sqrt3)\)，\(\bs b=(\sqrt3,1)\)，则 \(\bs a\) 与 \(\bs b\) 夹角的大小为\fillinblank{}.
\end{problem}

\begin{problem}
函数 \(f(x)=\dfrac{x}{x-1}\)（\(x\ge 2\)）的最大值为\fillinblank{}.
\end{problem}

\begin{problem}
某四棱柱的三视图如图所示，则该四棱柱的体积为\fillinblank{}.

\examfiguregroup[float=inline]{img/2016/beijing_liberal/q11_fig1.png}{%
    \bitmapinclude[max width=0.52\linewidth]{img/2016/beijing_liberal/q11_fig1.png}\hfill
    \bitmapinclude[max width=0.18\linewidth]{img/2016/beijing_liberal/q11_fig2.png}\hfill
    \bitmapinclude[max width=0.26\linewidth]{img/2016/beijing_liberal/q11_fig3.png}%
}
\end{problem}

\begin{problem}
已知双曲线 \(\dfrac{x^2}{a^2}-\dfrac{y^2}{b^2}=1\)（\(a>0\)，\(b>0\)）的一条渐近线为 \(2x+y=0\)，一个焦点为 \((\sqrt5,0)\)，则 \(a=\fillinblank{}\)，\(b=\fillinblank{}\).
\end{problem}

\begin{problem}
在 \(\triangle ABC\) 中，\(\angle A=\dfrac{2\pi}{3}\)，\(a=\sqrt3\,c\)，则 \(\dfrac{b}{c}=\fillinblank{}\).
\end{problem}

\begin{problem}
某网店统计了连续三天售出商品的种类情况：第一天售出 \(19\) 种商品，第二天售出 \(13\) 种商品，第三天售出 \(18\) 种商品；前两天都售出的商品有 \(3\) 种，后两天都售出的商品有 \(4\) 种，则该网店

\begin{enumerate}
\item 第一天售出但第二天未售出的商品有\fillinblank{} 种；
\item 这三天售出的商品最少有\fillinblank{} 种.
\end{enumerate}
\end{problem}

\section{解答题}

\begin{problem}
已知 \(\{a_n\}\) 是等差数列，\(\{b_n\}\) 是等比数列，且 \(b_2=3\)，\(b_3=9\)，\(a_1=b_1\)，\(a_{14}=b_4\).

\begin{enumerate}
\item 求 \(\{a_n\}\) 的通项公式；
\item 设 \(c_n=a_n+b_n\)，求数列 \(\{c_n\}\) 的前 \(n\) 项和.
\end{enumerate}
\end{problem}

\begin{problem}
已知函数 \(f(x)=2\sin\omega x\cos\omega x+\cos2\omega x\)（\(\omega>0\)）的最小正周期为 \(\pi\).

\begin{enumerate}
\item 求 \(\omega\) 的值；
\item 求 \(f(x)\) 的单调递增区间.
\end{enumerate}
\end{problem}

\begin{problem}
某市居民用水拟实行阶梯水价，每人每月用水量中不超过 \(w\) 立方米的部分按 \(4\) 元/立方米收费，超出 \(w\) 立方米的部分按 \(10\) 元/立方米收费，从该市随机调查了 \(10000\) 位居民，获得了他们某月的用水量数据，整理得到如图频率分布直方图：

\begin{enumerate}
\item 如果 \(w\) 为整数，那么根据此次调查，为使 \(80\%\) 以上居民在该月的用水价格为 \(4\) 元/立方米，\(w\) 至少定为多少？
\item 假设同组中的每个数据用该组区间的右端点值代替，当 \(w=3\) 时，估计该市居民该月的人均水费.
\end{enumerate}

\bitmapfigure[max width=0.56\linewidth]{img/2016/beijing_liberal/q17_fig1.png}
\end{problem}

\begin{problem}
如图，在四棱锥 \(P-ABCD\) 中，\(PC\perp\) 平面 \(ABCD\)，\(AB\parallel DC\)，\(DC\perp AC\).

\begin{enumerate}
\item 求证：\(DC\perp\) 平面 \(PAC\)；
\item 求证：平面 \(PAB\perp\) 平面 \(PAC\)；
\item 设点 \(E\) 为 \(AB\) 的中点，在棱 \(PB\) 上是否存在点 \(F\)，使得 \(PA\parallel\) 平面 \(CEF\)？说明理由.
\end{enumerate}

\bitmapfigure[max width=0.52\linewidth]{img/2016/beijing_liberal/q18_fig1.png}
\end{problem}

\begin{problem}
已知椭圆 \(C:\dfrac{x^2}{a^2}+\dfrac{y^2}{b^2}=1\) 过点 \(A(2,0)\)，\(B(0,1)\) 两点.

\begin{enumerate}
\item 求椭圆 \(C\) 的方程及离心率；
\item 设 \(P\) 为第三象限内一点且在椭圆 \(C\) 上，直线 \(PA\) 与 \(y\) 轴交于点 \(M\)，直线 \(PB\) 与 \(x\) 轴交于点 \(N\)，求证：四边形 \(ABNM\) 的面积为定值.
\end{enumerate}

\bitmapfigure[max width=0.56\linewidth]{img/2016/beijing_liberal/q19_fig1.png}
\end{problem}

\begin{problem}
设函数 \(f(x)=x^3+ax^2+bx+c\).

\begin{enumerate}
\item 求曲线 \(y=f(x)\) 在点 \((0,f(0))\) 处的切线方程；
\item 设 \(a=b=4\)，若函数 \(f(x)\) 有三个不同零点，求 \(c\) 的取值范围；
\item 求证：\(a^2-3b>0\) 是 \(f(x)\) 有三个不同零点的必要而不充分条件.
\end{enumerate}
\end{problem}
